% LOI DA SUA (batch 17): XUNG DOT SO MENH DE trong phu luc rieng. Phu luc co
% proposition RIENG (prop:subsid trong blocks_vi/08_subsidence) va remark rieng
% (blocks_vi/09_cost). \newtheorem dung counter rieng nen truoc day prop:subsid in ra
% 'Mệnh đề 1' — trong khi CUNG PDF nay viet 'bị mù theo Mệnh đề 1 của bài chính'
% (stub r@prop:blind=1) va S11 la 'Chứng minh Mệnh đề 1 và Mệnh đề 2'. Nhan
% 'Mệnh đề 1' tro HAI menh de khac nhau trong mot tai lieu.
% Bai chinh co 4 menh de (blind=1, coincide=2, noalarm=3, complete=4) va 1 ghi chu,
% nen dat counter tiep noi => prop:subsid=5, remark=2: DUNG BANG so ma ban gop
% (main_hujos_vi + supp_body_vi) da in ra. Idempotent trong ban gop vi counter da
% o 4/1 khi bat dau phu luc.
% Dat o DAY (khong phai supplementary_vi.tex) vi gate check_supp_vi_numbers doc
% supp_body_vi.tex cho ban VI va supplementary.tex cho ban EN; dat o file gate khong
% doc se lam EN thua 2 so -> gate bao lech gia.
\setcounter{proposition}{4}
\setcounter{remark}{1}

\section*{S1. Các lớp gian lận và mỗi tầng phòng thủ quan sát được gì}

\begin{minipage}{\textwidth}\footnotesize
\captionof{table}{Ba đại lượng làm giả được của \eqref{eq:three} trong bài chính và lớp gian lận
nào làm thay đổi từng đại lượng. ``đúng'' nghĩa là đại lượng đó được giữ nguyên, nên tầng
phòng thủ tương ứng bị mù theo Mệnh đề~\ref{prop:blind} của bài chính.}
% Sinh TU DONG tu tables/tab_threat.tex (ban tieng Anh, von sinh tu outputs/*.csv)
% boi experiments/make_tables_hujos_vi.py. KHONG sua tay.
% Nhan dich qua bang tra cuong che; so chi doi dinh dang sang kieu Viet Nam (dau phay thap phan).
% Da assert: cung so dong, cung so vach, cung tap chu so voi ban EN.

% Sinh TU DONG tu outputs/eval_theta_*.csv boi experiments/make_tables_anchored.py
% KHONG sua tay: moi con so deu truy nguon duoc ve CSV.

\fittocol{%
\begin{tabular}{@{}>{\raggedright\arraybackslash}p{2.85cm}>{\raggedright\arraybackslash}p{0.95cm}>{\raggedright\arraybackslash}p{1.05cm}>{\raggedright\arraybackslash}p{0.95cm}>{\raggedright\arraybackslash}p{2.20cm}@{}}
\toprule
Biến thể gian lận & Tổng mm & Từng độ sâu & Thời điểm & Làm giả cái gì \\
\midrule
Nông dân trung thực (điều kiện ruộng) & đúng & đúng & đúng & Mốc: bảy trục khác biệt thực tế đang bật \\
Thêm lệnh tưới & SAI & đúng & đúng & Thổi phồng hồ sơ hoạt động quản lý \\
Độ sâu khai không khả thi & SAI & SAI & đúng & Một phần ba số mục đặt bằng toàn bộ dung tích chứa \\
Bỏ sót lệnh thật & SAI & đúng & đúng & Xoá cửa sổ 96 h để tạo ra một pha khô \\
Gộp khai báo, 24 h & đúng & thay đổi & SAI & Cùng lượng nước, ít lần tưới hơn và lớn hơn \\
Gộp khai báo, 96 h & đúng & thay đổi & SAI & Cùng lượng nước, các lần tưới bám vào lưới thô \\
Dịch thời điểm (duy lý) & đúng & đúng & SAI & Leo đồi theo từng giờ để tối đa số pha khô được cấp tín chỉ \\
Bên gian lận duy lý & đúng & đúng & SAI & Chọn cửa sổ gộp làm tối đa số pha khô \\
\bottomrule
\end{tabular}
}

\end{minipage}

% ---------------------------------------------------------------- S2
\section*{S2. Biên độ sâu tưới chấp nhận được trên toàn $\Theta$}

\begin{minipage}{\textwidth}\footnotesize
\captionof{table}{Biên độ sâu tưới chấp nhận được $u_{\max}$ (\eqref{eq:umax} trong bài chính)
trên toàn dải của FAO-56 Bảng 2. Giá trị đơn lẻ $61$\,mm ban đầu chỉ là một dòng của bảng này.}
% Sinh TU DONG tu tables/tab_umax.tex (ban tieng Anh, von sinh tu outputs/*.csv)
% boi experiments/make_tables_hujos_vi.py. KHONG sua tay.
% Nhan dich qua bang tra cuong che; so chi doi dinh dang sang kieu Viet Nam (dau phay thap phan).
% Da assert: cung so dong, cung so vach, cung tap chu so voi ban EN.

% Sinh TU DONG tu outputs/eval_theta_*.csv boi experiments/make_tables_anchored.py
% KHONG sua tay: moi con so deu truy nguon duoc ve CSV.

\fittocol{%
\begin{tabular}{@{}>{\raggedright\arraybackslash}p{3.00cm}>{\raggedright\arraybackslash}p{1.25cm}>{\raggedright\arraybackslash}p{1.85cm}>{\raggedright\arraybackslash}p{1.85cm}@{}}
\toprule
Loại đất (FAO-56 Bảng 2) & Biên độ (mm) & $u_{\max}$ tại $w_0=0$ & $u_{\max}$ tại $w_0=0{,}90$ \\
\midrule
đất thịt pha cát & 55,0 & 49,0 & 54,5 \\
đất thịt & 75,0 & 65,0 & 72,5 \\
đất thịt pha limon & 90,0 & 77,0 & 86,0 \\
đất thịt pha sét limon & 70,0 & 61,0 & 68,0 \\
đất sét & 80,0 & 69,0 & 77,0 \\
\midrule
\textbf{Dải phải công bố} & --- & \textbf{49,0} & \textbf{86,0} \\
\bottomrule
\end{tabular}
}

\end{minipage}

% ---------------------------------------------------------------- S3
\section*{S3. Khả năng kích hoạt của bộ lọc điểm héo}

\begin{minipage}{\textwidth}\footnotesize
\captionof{table}{Khả năng kích hoạt của bộ lọc điểm héo trên $150$ cấu hình với nhật ký rỗng, tức
điều kiện thời tiết khô hạn nhất có được.}
% Sinh TU DONG tu tables/tab_f2_unreachable.tex (ban tieng Anh, von sinh tu outputs/*.csv)
% boi experiments/make_tables_hujos_vi.py. KHONG sua tay.
% Nhan dich qua bang tra cuong che; so chi doi dinh dang sang kieu Viet Nam (dau phay thap phan).
% Da assert: cung so dong, cung so vach, cung tap chu so voi ban EN.

% Sinh TU DONG tu outputs/eval_theta_*.csv boi experiments/make_tables_anchored.py
% KHONG sua tay: moi con so deu truy nguon duoc ve CSV.

\fittocol{%
\begin{tabular}{@{}>{\raggedright\arraybackslash}p{1.55cm}c>{\raggedright\arraybackslash}p{1.75cm}>{\raggedright\arraybackslash}p{0.85cm}>{\raggedright\arraybackslash}p{2.60cm}@{}}
\toprule
Trạm & Năm & Cạn kiệt cắt ngọn lớn nhất & $V_2$ kích hoạt & Diễn giải \\
\midrule
Cần Thơ & 2024 & 0,998 & 0 & không bao giờ xuống dưới điểm héo \\
Cần Thơ & 2025 & 0,988 & 0 & không bao giờ xuống dưới điểm héo \\
Sóc Trăng & 2024 & 0,999 & 0 & không bao giờ xuống dưới điểm héo \\
Sóc Trăng & 2025 & 0,992 & 0 & không bao giờ xuống dưới điểm héo \\
Cà Mau & 2024 & 0,991 & 0 & không bao giờ xuống dưới điểm héo \\
Cà Mau & 2025 & 0,980 & 0 & không bao giờ xuống dưới điểm héo \\
\bottomrule
\end{tabular}
}

\end{minipage}

% ---------------------------------------------------------------- S4
% S4: ban TIENG VIET cua paper2/blocks/08_subsidence.tex.
% Chi duoc \input boi paper2_hujos/supplementary_vi.tex.
% NGUYEN TAC: giu nguyen moi con so va moi hedge; thap phan kieu Viet ({,});
% tham chieu cheo toi bai chinh dung \eqref/\ref (so lay tu main_hujos_vi.aux
% boi experiments/make_supp_stubs_vi.py) thay vi go tay "Eq.~(16)" — ban EN go tay
% nen da lech so voi ban VI (eq:umax EN=16, VI=15).
\section*{S4. Tính bất biến đối với sụt lún}

Tuân thủ AWD thường được phát biểu bằng một mực nước ngầm hoặc độ sâu ngập tuyệt
đối, ví dụ $15$\,cm dưới mặt ruộng, tức ngưỡng tưới lại của nghiệm thức Safe-AWD
trong~\cite{balaine2019}. Ở một đồng bằng đang sụt lún, cao độ mặt ruộng $z_s(t)$ hạ
với tốc độ $\dot z$, nên một ngưỡng tuyệt đối đo theo một mốc cố định sẽ trôi.

\begin{proposition}[Độ trôi của một ngưỡng tuyệt đối]
\label{prop:subsid}
Trong một chu kỳ cấp tín chỉ dài $T_c$ năm, một ngưỡng tuyệt đối biểu diễn theo một
cao độ gốc cố định sẽ dịch đi $\dot z\,T_c$. Với $\dot z\in[1{,}1;2{,}5]$~cm/năm và
$T_c=5$~năm, độ dịch là $5{,}5$--$12{,}5$~cm, tức $37$--$83\%$ của một ngưỡng
$15$\,cm. Một tiêu chí viết theo tỷ lệ cạn kiệt \eqref{eq:depl} có độ trôi bằng
không, vì $FC$ và $WP$ là tính chất của đất quy về tầng rễ chứ không quy về một cao
độ tuyệt đối.
\end{proposition}

\begin{proof}
Tiêu chí tuyệt đối so một độ sâu đo được với một mốc không đổi; nếu mốc và mặt ruộng
rời nhau với tốc độ $\dot z$ thì ngưỡng hữu hiệu thay đổi $\dot z T_c$ sau $T_c$ năm.
Tiêu chí cạn kiệt so $FC-w_t$ với $(FC-WP)$, cả hai đều được định nghĩa bên trong tầng
rễ và tịnh tiến cứng cùng mặt ruộng, nên tỷ số của chúng bất biến đối với $z_s$.
\end{proof}

Các tốc độ sụt lún là số đo cho đồng bằng này trong~\cite{minderhoud2017}. Hệ quả thực
tiễn là một phương án neo theo tỷ lệ cạn kiệt không cần đo đạc lại các mốc ruộng trong
suốt chu kỳ cấp tín chỉ năm năm, còn phương án dùng độ sâu tuyệt đối thì cần.


% ---------------------------------------------------------------- S5
% S5: ban TIENG VIET cua paper2/blocks/09_cost.tex.
% Chi duoc \input boi paper2_hujos/supplementary_vi.tex.
% \label{sec:cost} GIU LAI vi ban EN van \ref no.
% Chu trong \underbrace dich theo dung ten cot cua Bang S5 (tables_vi/tab_baseline):
% phan cung / ket noi / nhan cong / tham dinh.
\section*{S5. Mô hình chi phí và điểm hòa vốn}
\label{sec:cost}

Giả sử một hợp tác xã phục vụ $A$ héc-ta với $N$ vụ mùa trong đời thiết bị, và chi phí
MRV mỗi héc-ta mỗi vụ phân tích thành
\begin{equation}
C=\underbrace{\frac{n_{\mathrm{dev}}\,p_{\mathrm{dev}}}{A\,N}}_{\text{phần cứng}}
+\underbrace{c_{\mathrm{conn}}}_{\text{kết nối}}
+\underbrace{h\,w}_{\text{nhân công}}
+\underbrace{\frac{c_{\mathrm{audit}}}{A}}_{\text{thẩm định}},
\label{eq:cost}
\end{equation}
trong đó $n_{\mathrm{dev}}$ và $p_{\mathrm{dev}}$ là số lượng và đơn giá thiết bị, $h$
là số giờ công mỗi héc-ta mỗi vụ và $w$ là tiền công nông thôn. Doanh thu mỗi héc-ta là
\begin{equation}
R=\kappa\,\rho ,
\label{eq:revenue}
\end{equation}
với $\kappa$ là số tín chỉ cấp cho mỗi héc-ta và $\rho$ là giá mỗi tín chỉ. Phương án
khả thi về kinh tế khi
\begin{equation}
C\ \le\ \beta R ,\qquad \beta\in(0,1),
\label{eq:breakeven}
\end{equation}
và $\beta=0{,}5$ là mức mà giám sát ngốn hết một nửa doanh thu tín chỉ. Thế
\eqref{eq:revenue} vào ta được lượng tín chỉ tối thiểu để hòa vốn
\begin{equation}
\kappa_{\min}=\frac{C}{\beta\rho}.
\label{eq:kappamin}
\end{equation}

Các tầng khác nhau ở $n_{\mathrm{dev}}$ và, quan trọng hơn, ở chỗ chúng quan sát được
đại lượng nào trong ba đại lượng làm giả được của \eqref{eq:three}. Tầng~0 là tầng vật
lý: $n_{\mathrm{dev}}=0$ ngoài bộ điều khiển, nhật ký được ghi trên điện thoại mà nông
dân đã có sẵn, và các bộ lọc ở Mục~\ref{sec:filters} chạy trên dữ liệu thời tiết công
khai với chi phí biên bằng không. Tầng~1 thêm một đồng hồ lưu lượng tại trạm bơm của hợp
tác xã --- $n_{\mathrm{dev}}=1$ cho toàn bộ diện tích $A$ chứ không phải mỗi thửa một
cái --- và quan sát được tổng lượng nước. Tầng~2 thêm một nút đo mực nước cho mỗi hợp
tác xã, neo lại bản sao số và khôi phục tính hiệu chuẩn mà một hệ MRV dựa trên đo đạc
dù sao cũng cần.

Không tầng nào trong ba tầng đó quan sát được dấu thời điểm, và theo
Mệnh đề~\ref{prop:blind} thì hai tầng đầu bất biến chính xác đối với phép biến đổi sinh
lợi. Vì vậy neo thiết bị ở Mục~\ref{sec:anchor} không phải là một tầng thứ tư tuỳ chọn
mà là điều kiện để bằng chứng được chấp nhận: nó do chính bộ điều khiển đóng máy bơm
sinh ra, nên chỉ thêm việc lưu trữ và truyền đi một bản ghi vốn đã tồn tại, và không
thêm $n_{\mathrm{dev}}$ nào.

Bảng~\ref{tab:baseline} so sánh phương án này với ba kiến trúc mà nó thay thế. Với
$\kappa\in[4,6]$ và $\rho\in[\$15,\$25]$~\cite{rmi2025rice,upstream2026price}, doanh thu
là $R\in[\$60,\$150]$ mỗi héc-ta mỗi vụ, nên $\kappa_{\min}$ tại $\beta=0{,}5$ tương ứng
với một ngân sách giám sát $30$--$75$~USD mỗi héc-ta. Toàn bộ ngăn xếp của bài tốn
$8{,}4$ USD mỗi héc-ta mỗi vụ so với ngân sách $30$--$75$ đó, tức thấp hơn trần từ
$3{,}6$ đến $8{,}9$ lần; vì vậy biên độ được phát biểu thành một dải chứ không phải một
bậc độ lớn. Nó cũng không phải dòng rẻ nhất của Bảng~\ref{tab:baseline}: viễn thám kèm
kiểm tra thực địa tốn $8{,}0$, tức rẻ hơn $0{,}4$. Hai phương án trả lời hai câu hỏi
khác nhau --- phương án rẻ hơn xác nhận rằng một pha khô đã xảy ra nhưng không giữ bản
ghi theo từng lệnh nào để đối chiếu với một nhật ký bị làm giả --- và chi phí tăng thêm
của neo so với một bộ điều khiển đã lắp sẵn là khoảng $1\%$ doanh thu ở mức thận trọng.
Bảng~\ref{tab:baseline} phân tích phép so sánh với ba kiến trúc mà nó thay thế; hạng mục
kết nối được tách riêng vì đó chính là thứ phân biệt một thiết kế mỗi thửa một cảm biến
với một gateway dùng chung.

\begin{remark}
Giá thiết bị trong Bảng~\ref{tab:baseline} là ước tính phục vụ lập kế hoạch, không phải
báo giá. Thứ tự giữa các tầng và bất đẳng thức hòa vốn \eqref{eq:breakeven} không phụ
thuộc vào chúng; chỉ có biên độ tuyệt đối là phụ thuộc. Cần thay các ước tính bằng báo
giá trước bất kỳ quyết định mua sắm nào.
\end{remark}


% RIENG khoi nay dung float [!ht] + FloatBarrier (Batch 5c): bang S4 rong gan
% sat \columnwidth, dat trong minipage thi \fittocol/\resizebox bao Overfull
% 28.45pt (ban goc dung float: 0 overfull). Float von giu caption+bang dinh
% nhau; FloatBarrier (placeins) chan khong cho bang troi qua S6.
\begin{table}[!ht]
\centering\footnotesize
\caption{Chi phí MRV mỗi héc-ta mỗi vụ cho bốn kiến trúc. Giám sát chiếm phần lớn chi phí
MRV, nên kiến trúc tránh được cảm biến chuyên dụng thắng ở phần nhân công chứ không phải
ở phần cứng.}\label{tab:baseline}
% Sinh TU DONG tu tables/tab_baseline.tex (ban tieng Anh, von sinh tu outputs/*.csv)
% boi experiments/make_tables_hujos_vi.py. KHONG sua tay.
% Nhan dich qua bang tra cuong che; so chi doi dinh dang sang kieu Viet Nam (dau phay thap phan).
% Da assert: cung so dong, cung so vach, cung tap chu so voi ban EN.

% Sinh TU DONG tu outputs/eval_theta_*.csv boi experiments/make_tables_anchored.py
% KHONG sua tay: moi con so deu truy nguon duoc ve CSV.

\fittocol{%
\begin{tabular}{@{}>{\raggedright\arraybackslash}p{3.00cm}>{\raggedright\arraybackslash}p{1.10cm}>{\raggedright\arraybackslash}p{1.30cm}>{\raggedright\arraybackslash}p{1.00cm}>{\raggedright\arraybackslash}p{0.85cm}>{\raggedright\arraybackslash}p{1.10cm}>{\raggedright\arraybackslash}p{4.00cm}@{}}
\toprule
Kiến trúc & Phần cứng & Kết nối & Nhân công & Thẩm định & Tổng \$/ha & Ghi chú \\
\midrule
Ống quan trắc thủ công kèm ảnh chụp & 0,5 & 0,0 & 8,0 & 4,0 & \textbf{12,5} & thực hành hiện hành theo Quyết định 4801 \\
Cảm biến mực nước mỗi thửa & 11,3 & 2,0 & 1,5 & 4,0 & \textbf{18,8} & chính xác, khó mở rộng cho nông hộ nhỏ \\
Viễn thám kèm kiểm tra thực địa & 0,0 & 0,0 & 2,0 & 6,0 & \textbf{8,0} & chu kỳ chụp 2 đến 4 ngày \\
Nhật ký được chứng thực, bài này & 3,6 & 1,5 & 0,3 & 3,0 & \textbf{8,4} & dấu thời gian của gateway là thứ làm nó được chấp nhận \\
\bottomrule
\end{tabular}
}

\end{table}
\FloatBarrier

% ---------------------------------------------------------------- S6
\section*{S6. Cấu hình thực nghiệm}

\begin{minipage}{\textwidth}\footnotesize
\captionof{table}{Cấu hình thực nghiệm. Mọi giá trị đều tái lập được từ mã nguồn đã công bố và dữ
liệu ERA5 công khai.}\label{tab:setup}
% Sinh TU DONG tu tables/tab_setup.tex (ban tieng Anh, von sinh tu outputs/*.csv)
% boi experiments/make_tables_hujos_vi.py. KHONG sua tay.
% Nhan dich qua bang tra cuong che; so chi doi dinh dang sang kieu Viet Nam (dau phay thap phan).
% Da assert: cung so dong, cung so vach, cung tap chu so voi ban EN.

% Sinh TU DONG tu outputs/eval_theta_*.csv boi experiments/make_tables_anchored.py
% KHONG sua tay: moi con so deu truy nguon duoc ve CSV.

\fittocol{%
\begin{tabular}{@{}>{\raggedright\arraybackslash}p{2.85cm}>{\raggedright\arraybackslash}p{5.60cm}@{}}
\toprule
Thành phần & Giá trị \\
\midrule
Thời tiết & ERA5 theo giờ qua kho lưu trữ Open-Meteo, không có giá trị thiếu \\
Trạm & 3: Cà Mau, Cần Thơ, Sóc Trăng \\
Năm khí hậu & 2024, 2025 (2024 khan nước, 2025 thừa nước) \\
Cửa sổ & từ ngày 1 tháng Ba đến ngày 29 tháng Năm (2\,160 h), mùa khô đồng bằng \\
Loại đất được lấy mẫu & 5 loại trải theo FAO-56 Bảng 2 \\
Trục khác biệt thực tế & 7: đất, giai đoạn $K_c$, $\eta$, mưa tại ruộng, thấm bờ, hành vi ghi chép, $w_0$ \\
Hiệu suất tưới $\eta$ & 0,60--0,85 (nước tới rễ / nước bơm) \\
Chế độ $K_c$ & cố định 1,05 (như đã công bố) và theo giai đoạn 1,05/1,20/0,90 \\
Nhiễu ghi chép & quên 5\%, làm tròn đến 5\,mm, trễ một ngày 35\%, sai độ sâu 15\% trong $\pm$35\% \\
Biến thể gian lận & 7 (ba ngây thơ, bốn đối kháng) \\
Tập bất định $|\Theta|$ & 30 = 5 loại đất $\times$ 2 chế độ $K_c$ $\times$ 3 hiệu suất \\
Trạng thái đầu $|w_0|$ & 5, trải cạn kiệt từ 0 đến 0,90 \\
Số lần phát lại mỗi nhật ký & 24 ($\Theta$ rút gọn) và 150 ($\Theta$ đầy đủ) \\
Số nhật ký được kiểm & 281 \\
\bottomrule
\end{tabular}
}

\end{minipage}

% ---------------------------------------------------------------- S7
\section*{S7. Hai chi tiết triển khai của phép giao bền vững}

Hai chi tiết triển khai của $V_{\mathrm{rob}}$ là quan trọng, và sai chi tiết nào cũng âm
thầm phá vỡ Mệnh đề~\ref{prop:noalarm} của bài chính. Một thế giới không sinh vi phạm nào
phải đóng góp \emph{tập rỗng} vào phép giao; bỏ qua các thế giới không kích hoạt là tính
$\bigcap$ chỉ trên các thế giới có kích hoạt, tức một lượng từ tồn tại mặc áo lượng từ phổ
dụng. Và $\mathcal{W}$ phải trải hết dải mà một khai báo còn có thể hợp lệ: với
$\mathcal{W}=\{0;0{,}25\}$ thì độ sâu chấp nhận được lớn nhất trong $\Theta$ là
$22{,}5$\,mm, nên mọi khai báo vượt mức đó đều ``có tội trong mọi thế giới'' và phép giao
thành rỗng nghĩa --- nó kích hoạt trên $102$ trong $150$ thế giới trong khi vẫn có $48$ thế
giới vô tội. Mở rộng $\mathcal{W}$ thành $\{0;0{,}25;0{,}5;0{,}75;0{,}9\}$ khôi phục đúng
ngữ nghĩa dự kiến.

% ---------------------------------------------------------------- S8
\section*{S8. Khảo sát đầy đủ các công trình giám sát trước đây}

Giám sát AWD được tiếp cận từ ba hướng. \emph{Gắn thiết bị}: cảm biến mực nước không dây
giá rẻ cho lúa AWD~\cite{delacruz2022} giải được bài toán đo lường nhưng nhân chi phí
thiết bị theo số thửa. \emph{Viễn thám}: nhận diện AWD ở đồng bằng sông Cửu
Long~\cite{lovell2019}, giám sát trạng thái nước bằng phân tách công suất tán xạ
ALOS-2~\cite{islam2026alose}, lập bản đồ học máy quy mô vùng cho các thực hành tiết kiệm
nước~\cite{mlrice2025}; phủ diện rộng với chi phí thấp nhưng hai đến bốn ngày mới chụp lại
nên không phân giải được pha khô $72$\,h vốn quyết định điều kiện cấp tín chỉ. \emph{Cấp
chương trình}: chương trình một triệu héc-ta lúa chất lượng cao, phát thải
thấp~\cite{qd1490} đến năm 2026 đã đưa hơn $350\,000$ héc-ta vào canh
tác~\cite{mae2026progress}, còn MRV được xác định là mắt xích thiếu của lúa phát thải thấp
ở Đông Nam Á~\cite{irri2026mrv}.

Về phía nông học, phân tích tổng hợp định lượng tác động của AWD lên nước và năng
suất~\cite{carrijo2017}; kinh tế học về việc chấp nhận~\cite{lampayan2015} và các rào cản
kỹ thuật, thể chế khi mở rộng~\cite{enriquez2021} giải thích vì sao giám sát gắn thiết bị
chưa lan rộng. Không hướng nào đánh giá \emph{độ tin cậy của chính bản ghi hành động}, vốn
là câu hỏi đặt ra ở đây.

Gateway cung cấp dấu thời điểm neo theo thiết bị chính là bộ điều khiển trong bản thảo chưa công bố của cùng nhóm tác giả~\cite{twingate2026}, nơi việc phát lệnh chạy trên kênh nhiễu chùm loại Gilbert--Elliott~\cite{gilbert1960,elliott1963};
chi phí biên của việc chứng thực do Mục~\ref{sec:anchor} của bài chính phân tích.

% ---------------------------------------------------------------- S9
\section*{S9. Bàn luận mở rộng}

\subsection*{S9.1 Vì sao công bố ngưỡng không phải là vấn đề}

Có thể nghĩ rằng nên giữ bí mật $\varepsilon$. Không nên, vì hai lý do. Kiểm toán viên phải
tái lập được phán quyết, và điều đó đòi hỏi tham số được công bố. Giữ bí mật cũng chẳng ích
gì: theo Mệnh đề~\ref{prop:blind}, tấn công dịch thời điểm giữ nguyên tổng lượng và từng độ
sâu, nên nó không cần biết $\varepsilon$ để nằm dưới giới hạn thấm --- nó chỉ cần tránh dịch
chuyển nước, điều mà nông dân nào cũng nhìn thấy. Lỗ hổng có tính cấu trúc, không phải hệ
quả của việc công bố.

\subsection*{S9.2 Khảo sát đầy đủ về viễn thám}

Radar vệ tinh là phương án giá rẻ hiển nhiên, nên thay vì tranh luận chúng tôi đo nó.
Bảng~\ref{tab:benchmark} của bài chính chấm mọi phương án bằng một đại lượng: trong $150$
pha khô đủ điều kiện cấp tín chỉ có thật trên $60$ ruộng mô phỏng, phương án đó xác nhận
được bao nhiêu. Một pha chỉ được tính nếu có ít nhất một lần quan sát rơi vào trong nó, vì
một pha $72$\,h nằm trọn giữa hai lần chụp thì không thể có bằng chứng nào --- nông dân mất
một tín chỉ ngay cả khi thực hành đúng. Sentinel-1 một vệ tinh, chu kỳ mười hai
ngày~\cite{clauss2017s1}, xác nhận $90/150$ ($60\%$); cả chòm A+B với chu kỳ sáu
ngày~\cite{hashemi2022s1} đạt $124/150$ ($83\%$) và vẫn bỏ sót một pha trong mỗi sáu pha dù
chi phí thiết bị bằng không. Một ống đo mực nước liên tục xác nhận cả $150$, mỗi thửa một
thiết bị~\cite{delacruz2022}. Bản ghi lệnh có chữ ký cũng xác nhận cả $150$ --- vì các pha
được định nghĩa bằng những khoảng giữa hai lần đóng mở mà bộ điều khiển vốn đã đóng dấu
thời điểm --- và không cần thêm thiết bị, vì chính bộ điều khiển đã mở máy bơm.

Điều bài chính chưa nói là hệ quả vận hành của hai con số đó. Phương án chụp ảnh cộng tự
khai báo không vô giá trị, nhưng vì nó không sinh ra chuỗi quan sát nào để đối chiếu nên
mọi xác nhận theo nó đều phải đến từ một chuyến đi thực địa: chi phí của nó nằm ở nhân công
chứ không nằm ở thiết bị, và đó là chỗ mà bản ghi lệnh có chữ ký thắng được.

% ---------------------------------------------------------------- S10
\section*{S10. Các lỗi của bộ khung thử nghiệm và chi tiết thiết kế}

\subsection*{S10.1 Hành vi ghi chép theo từng sự kiện chứ không theo từng giờ}

Việc làm tròn và trễ được áp theo từng \emph{sự kiện}, không theo từng giờ: áp mức làm tròn
$5$\,mm cho một sự kiện $40$\,mm rải trên tám giờ sẽ làm tròn mỗi mảnh $5$\,mm về không và
âm thầm xóa mất bản ghi, một ảo ảnh đã từng sinh ra các nhật ký rỗng trong phiên bản đầu
của bộ khung này.

\subsection*{S10.2 Hai lỗi bộ khung phải sửa}

Phải sửa hai lỗi của bộ khung trước khi các con số báo cáo có ý nghĩa, và cả hai đều được
ghi trong mã nguồn vì rất dễ bị đưa trở lại. Thứ nhất, trình sinh gian lận không tìm được
nước đi cải thiện sẽ trả lại nhật ký trung thực \emph{ban đầu}, nên đếm vô điều kiện sẽ gộp
cả những ca không hề có tấn công vào số ``tấn công bị bắt''; vì vậy mọi dòng đối kháng mang
cờ \texttt{changed} và mọi thống kê tấn công đều tính có điều kiện trên cờ đó. Thứ hai, mỗi
biến thể phải so với mốc của chính nó: nhật ký trung thực điều kiện giấy do một lần gọi
riêng sinh ra, và việc chấm nó theo bản ghi điều kiện ruộng của một lần hiện thực hóa đất
khác đã ``bắt'' sai $56$ trong $60$ nhật ký trung thực.

\subsection*{S10.3 Tổng lượng là chỉ báo chẩn đoán, không phải mục tiêu}

Một cấu hình tối đa hóa tổng lượng khai báo không tự động sinh lợi: một bên gian lận ở phiên bản sớm thổi phồng khai báo thêm $344$\,mm đã kéo $n_{\mathrm{dry}}$ về không và sẽ chẳng kiếm
được gì, nên tổng lượng chỉ được báo cáo như một chỉ báo chẩn đoán.

\subsection*{S10.4 Ảo ảnh tự nhất quán}

Tầng được kiểm toán từng sinh các quỹ đạo ``thật'' bằng \texttt{SoilProfile} mặc định, ERA5
không nhiễu, tưới là xung một giờ đúng bằng $\{20,40\}$\,mm, ngưỡng đúng bằng $0{,}80$, và
một nông dân luôn ghi đúng độ sâu vào đúng giờ. Vì vậy cả $0/60$ ca kết tội oan lẫn $60/60$
ca phát hiện đều chỉ đo tính tự nhất quán: bộ phát hiện đang chấm chính trình sinh đã tạo ra
các quỹ đạo ``thật'' mà nó kiểm toán. Trong thiết kế hiện tại, thửa ruộng do một tiến trình
riêng biệt sinh ra, khác bản sao số theo bảy trục liệt kê ở Bảng~S5.

% ---------------------------------------------------------------- S11
\section*{S11. Chứng minh Mệnh đề 1 và Mệnh đề 2}

\subsection*{Chứng minh Mệnh đề 1 về sự mù chính xác của các tầng phòng thủ tổng hợp}

$U(L)=\sum_i u_i$ là một hàm đối xứng của đa tập hợp, nên nó bất biến đối với mọi phép gán
lại độ sâu cho các giờ; $M$ là tính chất của thửa ruộng chứ không phải của nhật ký; do đó
$G=M-U$ không đổi. Tương tự, giới hạn thấm (\eqref{eq:v5} của bài chính) được đánh giá
theo từng sự kiện và chỉ phụ thuộc $(u_i,\tau_i)$, không bao giờ phụ thuộc $t_i$ hay
$\hat w$, nên số lần vi phạm là một hàm chỉ của đa tập hợp.

\subsection*{Chứng minh Mệnh đề 2 về lợi ích tín chỉ và sự lọt lưới là cùng một hành động}

Phần (i) theo định nghĩa: tập $\mathcal{A}$ không đòi hỏi khai báo nào bên trong khoảng, nên
bỏ một khai báo sẽ kéo dài hoặc tạo ra một pha đạt điều kiện. Với phần (ii), hoãn khai báo
đến $t+\delta$ nghĩa là thùng chứa tiếp tục mất nước thêm $\delta$ giờ mà không có nước vào,
nên $\hat w_{t+\delta}^{-}\le\hat w_t^{-}$ và $c_{t+\delta}\ge c_t$; cận trên của độ tăng phần
thiếu hụt suy ra vì lượng mất mỗi giờ ít nhất bằng trung bình $\overline{\Delta}$ cho tới khi
đạt sức chứa ẩm đồng ruộng, sau đó $c$ bão hòa tại $FC-\hat w$. Vì bộ lọc vượt sức chứa chỉ
kích hoạt khi $u_i>c_{t_i}+\varepsilon$ và $u_i$ không đổi, một $c_{t_i}$ lớn hơn chỉ có thể
làm giảm tập vi phạm.

% ---------------------------------------------------------------- S12
\section*{S12. Chứng minh Mệnh đề 4 về tính đầy đủ của neo}

Vế thứ nhất chính là định nghĩa của phán quyết neo: bất kỳ sự bất đồng nào giữa các tập giờ,
hoặc giữa độ sâu tại một giờ chung, đều kích hoạt phán quyết giả mạo. Với vế thứ hai, tổng
lượng, từng độ sâu riêng lẻ và các dấu thời điểm đều là hàm của tập giờ và độ sâu theo giờ,
nên sự trùng khớp trên hai thứ đó cố định cả ba, và không có sự giả mạo nào đối với ba đại
lượng đã xảy ra.

% ---------------------------------------------------------------- S13
\section*{S13. Toàn bộ dẫn xuất của bộ lọc điểm héo}

Tầng này từng gồm một bộ lọc vật lý thứ ba, đánh dấu việc cạn kiệt xuống dưới điểm héo là
bất khả thi về vật lý. Nó vô hiệu theo hai nghĩa. Ngưỡng đặt ở $1{,}02$ trong khi chỉ báo
được báo cáo $\hat d_t=\min(1,\max(0,\cdot))$ bị cắt về $[0,1]$, nên điều kiện không thể
thỏa mãn ngay từ cách dựng: $0$ lần kích hoạt trong $281$ nhật ký. Chúng tôi sửa bằng cách
đánh giá độ cạn kiệt chưa cắt ngọn, rồi đo xem bộ lọc đã sửa có kích hoạt được không. Trên
ba trạm, hai năm khí hậu, năm loại đất và năm trạng thái đầu --- tức $150$ cấu hình --- với
nhật ký rỗng, nghĩa là thời tiết khô hạn nhất có thể, độ cạn kiệt chưa cắt ngọn đạt đỉnh
$0{,}9987$ (Bảng~S3): ruộng lúa vụ khô ở đồng bằng không chạm tới điểm héo. Vậy bộ lọc này
là thiết bị \emph{an toàn} bảo đảm không trạng thái vô lý nào được chấp nhận, không phải
bằng chứng về năng lực phát hiện, và tầng này phải được mô tả là có hai bộ lọc vật lý hoạt
động và một bộ lọc trơ.

% ---------------------------------------------------------------- S14
\section*{S14. Hai hiệu chỉnh đối với biên độ sâu tưới chấp nhận được}

Thứ nhất, \eqref{eq:umax} của bài chính dùng $\max$ chứ không dùng một hiệu số: ruộng khô ban
đầu ($w_0=0{,}90$) chấp nhận những khai báo \emph{lớn hơn} ruộng ướt, còn biểu thức ban đầu
$(d^{*}-w_0)\,\mathrm{span}+\varepsilon$ cho ra $43{,}5$\,mm với $\mathrm{span}=70$\,mm và
$w_0=0{,}25$ --- một giá trị không tương ứng trạng thái vật lý nào, vì phần thiếu hụt hữu hiệu
tại giờ đó hoặc là $0{,}25\,\mathrm{span}$ (lúc đầu) hoặc là $d^{*}\mathrm{span}$ (khi ruộng đã
khô). Thứ hai, biên này áp theo \emph{từng sự kiện} chứ không theo từng vụ: thùng chứa rút
liên tục qua thấm sâu và bốc thoát hơi nước, nên tổng khai báo cả vụ không bị chặn bởi
$u_{\max}$ và một bên gian lận phát nhiều khai báo dưới ngưỡng có thể thổi phồng tổng lượng không
giới hạn.

% ---------------------------------------------------------------- S15
\section*{S15. Vì sao mốc ghép cặp mới là mốc trung thực}

Các biến thể giả mạo chỉ được sinh ở những quỹ đạo mang đủ lệnh để thao tác, tức từ ba lệnh
trở lên, và chúng vốn đã khô hơn trung bình: mốc trung thực của chúng mang $4{,}87$ pha khô
cấp được tín chỉ so với $3{,}23$ của toàn bộ nhật ký trung thực. Vì vậy so sánh không ghép cặp
với mốc rộng hơn sẽ quy phần khô hạn vốn có cho cuộc tấn công và thổi phồng mức tăng lên xấp xỉ
gấp đôi. Mọi con số trong Bảng~\ref{tab:credit} của bài chính đều tính trên tập con ghép cặp
--- cùng trạm, năm, hạt giống và loại đất --- có điều kiện trên cờ \texttt{changed}.

% ---------------------------------------------------------------- S16
\section*{S16. Chín loại nhật ký}

Ba loại là giả mạo ngây thơ: thêm sự kiện vào ngày mưa, thổi phồng một phần ba số mục lên bằng
toàn bộ biên độ chứa, và xóa một cửa sổ khai báo dài $96$\,h. Bốn loại là đối kháng và, khác các
loại ngây thơ, được dựng bằng bộ tối ưu chứ không theo công thức cố định: gộp khai báo về lưới
$24$\,h và $96$\,h; bên gian lận duy lý chọn cửa sổ gộp làm tối đa $n_{\mathrm{dry}}$ với điều kiện
giữ nguyên tổng lượng và tôn trọng giới hạn thấm; và phép leo đồi dịch thời điểm, tìm độ lệch
nguyên trong $[-36,36]$ và chỉ nhận một nước đi nếu tổng lượng giữ nguyên từng bít, mọi độ sâu
tưới nằm dưới giới hạn thấm, và $n_{\mathrm{dry}}$ tăng thực sự. Hai loại trung thực còn lại --- điều
kiện giấy và điều kiện ruộng --- là các mốc so sánh.

% ---------------------------------------------------------------- S17
\section*{S17. Cái giá của tính bền vững và sự nhất quán giữa các tập bất định}

\subsection*{S17.1 Mức mất mát năng lực phát hiện đo được}

Thay mô hình điểm bằng phép giao trên $\Theta$ không miễn phí, và mức mất mát phải được báo
cáo. Trên chín nhật ký gian lận bỏ sót thật sự bị sửa, mô hình điểm bắt được $9/9$ còn quy
tắc bền vững bắt $4/9$ ($44\%$); với $\Theta$ đầy đủ là $4/4$ và $2/4$ ($50\%$). Cùng một sự
đánh đổi chạy cùng hướng với ngưỡng tỉ lệ và vì cùng lý do: tầng phòng thủ từ chối kết tội
chỉ dựa trên một cấu hình mô hình thì cũng từ chối một số ca dương tính thật. Giả mạo ngây
thơ vẫn ở $100\%$ theo cả hai. Chúng tôi coi đây là phía đúng của sự đánh đổi, vì một ca kết
tội oan rơi xuống nông dân không có cách kháng cáo còn một ca bỏ lọt rơi xuống cơ quan xác
minh vốn vẫn lấy mẫu, nhưng đó là mất mát năng lực phát hiện chứ không phải cải thiện
Pareto.

\subsection*{S17.2 Tập rút gọn và tập đầy đủ cho cùng kết quả}

Mọi kết luận đều đúng với cả $\Theta$ rút gọn, gồm $24$ lần phát lại mỗi nhật ký, $281$ nhật ký
và $10$ hạt giống ngẫu nhiên, lẫn $\Theta$ đầy đủ, gồm $150$ lần phát lại, $85$ nhật ký và $3$
hạt giống ngẫu nhiên: kết tội oan $0/60$ và $0/18$, tính đầy đủ của neo $116/116$ và $37/37$,
lọt lưới dịch thời điểm ở tầng bền vững $8/9$ và $3/3$, thổi phồng tín chỉ ghép cặp $+33\%$
($t=5{,}29$) và $+38\%$ ($t=3{,}46$). Tập rút gọn được dùng cho các lần chạy lớn vì nó chứa cực
trị của mọi trục, tức nơi một thế giới vô tội dễ tồn tại nhất; tập đầy đủ xác nhận không cấu
hình nằm bên trong nào làm thay đổi phán quyết.

% ---------------------------------------------------------------- S18
\section*{S18. Vì sao neo không đưa việc kết tội oan trở lại}

Mệnh đề~\ref{prop:complete} của bài chính mạnh tới mức thoạt trông có vẻ nguy hiểm: mọi sai
lệch đều bị coi là giả mạo, kể cả quên ghi và làm tròn mà bài chính đã chỉ ra là rất phổ biến.
Cách hóa giải là neo không đối chiếu nhật ký với một \emph{mô hình} mà đối chiếu một bản chép
với \emph{nguồn}. Nông dân trung thực quên một mục không bị kết tội gian lận, vì hệ thống lưu
trữ hồ sơ là gateway chứ không phải trí nhớ của nông dân: khai báo đó đơn giản được bổ sung
từ $\mathcal{G}$. Về vận hành, neo biến nhật ký từ \emph{bằng chứng mà nông dân phải ghi cho
đúng} thành \emph{một góc nhìn của bản ghi thiết bị đã giữ sẵn}. Khi đó kết tội oan đòi hỏi
gateway bất đồng với chính nó, vốn là câu hỏi toàn vẹn thiết bị chứ không phải hành vi, và đó
là lý do bảo đảm này giữ ở $0/60$ trong khi các tầng dựa trên hành vi giữ ở $26{,}7\%$ và
$5{,}0\%$.

% ---------------------------------------------------------------- S19
\section*{S19. Cuộc tấn công lấy đi thứ gì}

Cuộc tấn công không lấy đi nước. Tổng lượng nước bơm không đổi, nên nông dân trả vẫn chừng
đó và khí quyển vẫn thấy chừng đó mê-tan. Thứ bị lấy đi là \emph{lịch sử}: tín chỉ được cấp
cho một mẫu pha khô đã không xảy ra. Theo VM0051, đơn vị cấp tín chỉ chính xác là mẫu đó, nên
thiệt hại là có thật dù không có nguồn lực vật lý nào dịch chuyển.

% ---------------------------------------------------------------- S20
\section*{S20. Toàn bộ chi tiết của thiết kế thực nghiệm}

\subsection*{S20.1 Gỡ bỏ tính vòng tròn, theo từng trục}

Thửa ruộng do một tiến trình riêng biệt sinh ra, khác bản sao số theo bảy trục, mỗi trục đều có
nguồn đã công bố (Bảng~S5): thành phần cơ giới của đất lấy từ năm cấp đất FAO-56 Bảng 2 có mặt ở
đồng bằng; $K_c$ phân theo giai đoạn đúng Bảng 12 chứ không cố định ở $1{,}05$; hiệu suất tưới
$\eta\in[0{,}60;0{,}85]$, nên nông dân khai lượng nước \emph{đã bơm} trong khi chỉ $\eta u$ tới
được tầng rễ; mưa tại ruộng bằng $0{,}7$--$1{,}4$ lần ô lưới cộng nhiễu Gauss; thấm bờ
$0$--$0{,}02$\,mm\,h$^{-1}$; hành vi ghi chép với $5\%$ quên ghi, làm tròn đến $5$\,mm, $35\%$
khả năng trễ một ngày và $15\%$ sai độ sâu tưới trong $\pm35\%$; và độ cạn kiệt ban đầu $w_0$ lấy từ
$[0;0{,}45]$ thay vì giả định là sức chứa ẩm đồng ruộng. Việc làm tròn và trễ áp theo từng
\emph{sự kiện} chứ không theo từng giờ (S10.1).

\subsection*{S20.2 Phần còn lại được của tuyên bố ban đầu}

Tiền đề thúc đẩy công trình này vốn hấp dẫn: một lệnh tưới tự thân đã là sự kiện số, nên bằng
chứng MRV lẽ ra gần như miễn phí. Tiền đề đó vẫn đứng vững nhưng không ở dạng tuyên bố ban
đầu: nó đứng vững vì bộ điều khiển phát lệnh cũng có thể ký lệnh đó, và chính bản ghi có chữ
ký làm cho bằng chứng được chấp nhận --- chứ không phải vì một mô hình thủy văn phát hiện được
lời khai dối. Ở chỗ bản thảo ban đầu nói vật lý \emph{chứng thực} nhật ký, phát biểu đúng là
vật lý \emph{ràng buộc} nhật ký: nó loại trừ những độ sâu không loại đất nào chứa được và
những trạng thái không thửa ruộng nào đạt tới --- điều đó đáng có --- và nó im lặng về đúng
phép biến đổi duy nhất sinh lợi. Giữ bí mật $\varepsilon$ cũng không bịt được lỗ hổng: cuộc tấn
công không cần biết ngưỡng, và dù sao kiểm toán viên cũng phải tái lập được phán quyết
(S9.1).

% ---------------------------------------------------------------- S21
\section*{S21. Bảy giới hạn của bài, phát biểu đầy đủ}

\begin{enumerate}
\item \emph{Nông dân mô phỏng.} Các hành vi giả mạo được tiêm vào những quỹ đạo tham chiếu mô
phỏng. Cần một thử nghiệm thực địa để đo phản ứng hành vi, cụ thể là liệu nông dân có thay
đổi việc bơm \emph{thực tế} của họ một khi biết rằng neo tồn tại hay không --- một kiểu tấn
công mà mô hình của chúng tôi không biểu diễn được, vì nó thay đổi $M$ chứ không thay đổi $L$.
\item \emph{Neo giả định một gateway trung thực.} Mệnh đề~\ref{prop:complete} của bài chính là
tính đầy đủ đối với việc giả mạo nhật ký, không phải đối với một bộ điều khiển đã bị chiếm
quyền. Toàn vẹn thiết bị, việc giữ khóa và chống phát lại được giả định từ công trình đồng
hành và không được thiết lập lại ở đây.
\item \emph{Không áp dụng cho tưới thủ công.} Nơi một cửa cống được mở bằng tay thì không có
bản ghi lệnh, nên neo không khả dụng. Điều này loại trừ một phần đáng kể số nông hộ ở đồng
bằng và giới hạn quần thể mà phương án có thể phục vụ.
\item \emph{Cuộc tấn công chỉ được minh họa trong năm khan nước.} Trong năm thừa nước 2025,
nông dân cần tưới ít đến mức phần lớn nhật ký mô phỏng rơi xuống dưới mức tối thiểu ba lệnh,
mức mà một tấn công thời điểm mới có thể dựng được: trong $37$ nhật ký thật sự bị sửa ở
$\Theta$ đầy đủ thì cả $37$ đều thuộc năm 2024. Chính vách ngăn giữa hai năm đó đã là một
phát hiện --- rủi ro gian lận MRV tập trung ở nơi khan nước --- nhưng nó cũng có nghĩa là các
kết quả tấn công chưa phải bằng chứng về một mùa mưa. Ở $\Theta$ rút gọn có mười hai nhật ký
bị sửa thuộc năm 2025 và chúng hành xử giống nhau, điều đó gợi ý chứ chưa kết luận được.
\item \emph{$\Theta$ là một phán đoán chuyên môn.} Năm loại đất, hai chế độ $K_c$, ba mức hiệu
suất và năm trạng thái đầu bao được đồng bằng nhưng không vét hết nó. Một thửa ruộng nằm ngoài
$\Theta$ sẽ mất bảo đảm của Mệnh đề~\ref{prop:noalarm} trong bài chính, và đó chính xác là
tình huống mà ngưỡng tỉ lệ cùng phán quyết chờ xác minh tồn tại để hấp thụ.
\item \emph{Giá thiết bị là ước tính phục vụ lập kế hoạch, không phải báo giá.} Thứ tự giữa các
tầng và bất đẳng thức hòa vốn không bị ảnh hưởng; biên độ tuyệt đối thì có.
\item \emph{Không định lượng phát thải.} Tầng này sinh ra dữ liệu hoạt động. Việc ghép nó với
một mô hình mê-tan và các hệ số phát thải được hiệu chuẩn tại chỗ là một công trình riêng,
cũng như giả định của bảng chi phí rằng tín chỉ tăng theo $n_{\mathrm{dry}}$.
\end{enumerate}
